Table~\ref{tab:main} gives the main comparison and Fig.~\ref{fig:curves} the learning curves (mean over seeds of the 1,000-run means). Bold and underline in the generated tables follow each metric's declared direction (lower is better) and are therefore not meaningful for the signed bias columns; we read the signed numbers directly. Unless stated, $p$ is the Welch $p$-value over 5 seeds with Double Q-learning as reference.

\begin{table*}[t]\centering
\caption{Main results (mean $\pm$ std over 5 seeds; each seed is the mean of 1{,}000 independent runs). bias\_first: start-state estimated-minus-true, mean over the first 20 episodes; bias\_final and qbias\_final: last 20 episodes; subopt\_all: suboptimal-action fraction over all episodes; greedy\_err\_final: greedy-policy error at the end. On the trap MDP bias is clipped at zero from below (Sec.~\ref{sec:method}).}\label{tab:main}
\resizebox{\textwidth}{!}{\begin{tabular}{llcccccc}
\toprule
Method & Task & bias\_first $\downarrow$ & bias\_final $\downarrow$ & qbias\_final $\downarrow$ & subopt\_all $\downarrow$ & greedy\_err\_final $\downarrow$ & $n$ \\
\midrule
Q-learning & maxbias & 0.038 $\pm$ 0.001 & 0.008 $\pm$ 0.002 & 0.075 $\pm$ 0.002 & 0.380 $\pm$ 0.005 & 0.083 $\pm$ 0.013 & 5 \\
Weighted Double Q-learning & maxbias & \underline{0.003 $\pm$ 0.000} & 0.001 $\pm$ 0.000 & \underline{0.073 $\pm$ 0.000} & \underline{0.127 $\pm$ 0.002} & \underline{0.020 $\pm$ 0.004} & 5 \\
Maxmin Q-learning & maxbias & \textbf{0.000 $\pm$ 0.000} & \textbf{0.001 $\pm$ 0.000} & 0.080 $\pm$ 0.000 & 0.197 $\pm$ 0.002 & 0.024 $\pm$ 0.002 & 5 \\
\textbf{Double Q-learning} & maxbias & 0.003 $\pm$ 0.000 & \underline{0.001 $\pm$ 0.000} & \textbf{0.061 $\pm$ 0.001} & \textbf{0.119 $\pm$ 0.002} & \textbf{0.017 $\pm$ 0.002} & 5 \\
\midrule
Q-learning & random20 & -2.960 $\pm$ 0.022 & -0.150 $\pm$ 0.011 & -0.444 $\pm$ 0.005 & \textbf{0.255 $\pm$ 0.002} & \textbf{0.173 $\pm$ 0.016} & 5 \\
Weighted Double Q-learning & random20 & \underline{-3.283 $\pm$ 0.016} & -0.360 $\pm$ 0.003 & -0.839 $\pm$ 0.005 & 0.313 $\pm$ 0.003 & 0.257 $\pm$ 0.010 & 5 \\
Maxmin Q-learning & random20 & \textbf{-3.428 $\pm$ 0.021} & \textbf{-0.721 $\pm$ 0.018} & \textbf{-1.175 $\pm$ 0.009} & \underline{0.282 $\pm$ 0.002} & \underline{0.228 $\pm$ 0.026} & 5 \\
\textbf{Double Q-learning} & random20 & -3.281 $\pm$ 0.017 & \underline{-0.364 $\pm$ 0.007} & \underline{-0.840 $\pm$ 0.005} & 0.312 $\pm$ 0.001 & 0.254 $\pm$ 0.011 & 5 \\
\bottomrule
\end{tabular}
}
\end{table*}

\begin{figure}[t]\centering\includegraphics[width=\columnwidth]{curves_main.pdf}
\caption{Top: start-state estimated-minus-true value; bottom: suboptimal-action fraction per episode. Left: trap MDP; right: random20 (y-axis clipped at $-4$ to $0.5$). Orange: Q-learning, green: Double, purple: Weighted Double, grey: Maxmin; Double and Weighted Double nearly coincide.}\label{fig:curves}
\end{figure}

\paragraph{H1 (trap MDP, Q-bias): supported in direction, not in the strong form.} The registered test favours Double: final $Q$-bias \rhval{main/double-q-learning/maxbias/qbias_final/mean} versus \rhval{main/q-learning/maxbias/qbias_final/mean} for Q-learning (difference \rhval{cmp/main/q-learning/maxbias/qbias_final/delta}, $p=$\rhval{cmp/main/q-learning/maxbias/qbias_final/welch_p}). But both values are positive, i.e.\ both agents still rate the left action higher than its true value $-0.1$ at the end of 300 episodes. A likely reason is the zero initialisation: the true value is below the initial value, and the left action is taken only about $\epsilon/2$ of the time late in training. The warm-start ablation (Sec.~\ref{sec:ablations}) removes this transient and gives a cleaner contrast.

\paragraph{H2 (trap MDP, actions): supported.} The suboptimal-action fraction is \rhval{main/double-q-learning/maxbias/subopt_all/mean} for Double versus \rhval{main/q-learning/maxbias/subopt_all/mean} for Q-learning ($p=$\rhval{cmp/main/q-learning/maxbias/subopt_all/welch_p}); in the first 20 episodes it is \rhval{main/double-q-learning/maxbias/subopt_first/mean} versus \rhval{main/q-learning/maxbias/subopt_first/mean}. The peak start-state overestimate, which is the textbook signature, is \rhval{main/q-learning/maxbias/bias_peak/mean} for Q-learning and \rhval{main/double-q-learning/maxbias/bias_peak/mean} for Double ($p=$\rhval{cmp/main/q-learning/maxbias/bias_peak/welch_p}). Weighted Double behaves like Double (peak \rhval{main/weighted-double-q-learning/maxbias/bias_peak/mean}). Maxmin keeps the bias small (\rhval{main/maxmin-q-learning/maxbias/bias_peak/mean}) but takes more suboptimal actions than Double (\rhval{main/maxmin-q-learning/maxbias/subopt_all/mean}): a possible cause is that two estimators with a minimum are slow to rank the actions at B, which we did not isolate.

\paragraph{H3 (random MDP): first part refuted, second part supported.} Q-learning's final start-state bias is \rhval{main/q-learning/random20/bias_final/mean}, i.e.\ negative, not positive, and Double Q-learning's is lower still (\rhval{main/double-q-learning/random20/bias_final/mean}; difference \rhval{cmp/main/q-learning/random20/bias_final/delta}, $p=$\rhval{cmp/main/q-learning/random20/bias_final/welch_p}). The all-pair $Q$-bias shows the same ordering (Q-learning \rhval{main/q-learning/random20/qbias_final/mean}, Double \rhval{main/double-q-learning/random20/qbias_final/mean}). Double does not take fewer suboptimal actions: \rhval{main/double-q-learning/random20/subopt_all/mean} versus \rhval{main/q-learning/random20/subopt_all/mean} for Q-learning ($p=$\rhval{cmp/main/q-learning/random20/subopt_all/welch_p}), and its final greedy error is higher (\rhval{main/double-q-learning/random20/greedy_err_final/mean} versus \rhval{main/q-learning/random20/greedy_err_final/mean}). Maxmin has the most negative bias (\rhval{main/maxmin-q-learning/random20/bias_final/mean}) but a lower suboptimal fraction than Double (\rhval{main/maxmin-q-learning/random20/subopt_all/mean}). Weighted Double is indistinguishable from Double here (bias \rhval{main/weighted-double-q-learning/random20/bias_final/mean}, suboptimal fraction \rhval{main/weighted-double-q-learning/random20/subopt_all/mean}). On this MDP family and at this budget, removing the positive bias therefore did not help; random20 is a regime where Q-learning is the best of the four on action quality, a failure case for double estimators.

\paragraph{What the numbers do not show.} The main table shows that the start-state estimates are below the true values for every method on random20, but it does not say why: at the end of training the zero initialisation, the per-table update rate and adaptive sampling all push the same way. The ablations below address each separately, and none of them is a proof of mechanism. The tests use $n=5$ seed means with very small spread, so the $p$-values are small for almost any difference; effect sizes in \texttt{rh compare} are large (e.g.\ Cohen's $d=$\rhval{cmp/main/q-learning/random20/subopt_all/cohen_d} for the suboptimal fraction on random20), and statistical significance here says only that the difference is reproducible across seeds of this MDP family.
